% !TEX TS-program = pdflatex
% !TeX program = pdflatex
% !TEX encoding = UTF-8
% !TEX spellcheck = en_US

\documentclass[xcolor=table]{beamer}

\usepackage{../extra/beamer/karimnlp}
\input{options}

\subtitle[11- Some applications]{Chapter 11\\Some applications} 

\changegraphpath{../img/app/}

\begin{document}

\begin{frame}
	\frametitle{\inserttitle}
	\framesubtitle{\insertshortsubtitle: Introduction}
	
	\begin{itemize}
		\item NLP applications can be categorized according to several criteria:
		\begin{itemize}
			\item Input modality: text or speech
			\item Interactivity: interactive vs. non-interactive systems
			\item Output type: structured labels (classes) or generated text
		\end{itemize}
		\item Based on these criteria, we distinguish four main categories:
		\begin{itemize}
			\item Text transformation
			\item Interactive systems
			\item Text classification
			\item Speech processing
		\end{itemize}
	\end{itemize}
	
\end{frame}


\begin{frame}
	\frametitle{\inserttitle}
	\framesubtitle{\insertshortsubtitle: Plan}

	\begin{multicols}{2}
	%	\small
	\tableofcontents
	\end{multicols}

\end{frame}

%===================================================================================
\section{Transformation}
%===================================================================================

\begin{frame}
	\frametitle{\insertshortsubtitle}
	\framesubtitle{\insertsection}
	
	\begin{itemize}
		\item Tasks where an input text is transformed into another text while preserving or compressing its meaning.
		\item \textbf{Input}: textual data
		\item \textbf{Output}: textual data
		\item \textbf{Applications:}
		\begin{itemize}
			\item Machine translation (MT)
			\item Automatic text summarization (ATS)
		\end{itemize}
	\end{itemize}

\end{frame}

\subsection{Machine translation}

\begin{frame}
	\frametitle{\insertshortsubtitle: \insertsection}
	\framesubtitle{\insertsubsection: Approaches}
	
	\hgraphpage{MT-classif.pdf}
	
\end{frame}

\begin{frame}
	\frametitle{\insertshortsubtitle: \insertsection}
	\framesubtitle{\insertsubsection: Direct}
	
	\begin{itemize}
		\item \textbf{Algorithm}:
		\begin{itemize}
			\item The source text $S$ is processed as a sequence of words.
			\item Each word $S_i$ is translated into a corresponding target word $T_i$ using a bilingual dictionary.
			\item The resulting sequence is then reordered to conform to the target language syntax (e.g., \expword{SVO $\rightarrow$ VSO, Adj + N $\rightarrow$ N + Adj}).
		\end{itemize}
		\item \textbf{Preconditions}:
		\begin{itemize}
			\item The source and target languages must be typologically similar (e.g., similar grammatical structures);
			\item Availability of a high-quality bilingual dictionary;
			\item Basic morphological processing tools.
		\end{itemize}
		\item \textbf{Level}: Primarily lexical (with limited syntactic adjustments).
		\item \textbf{Examples}: Early rule-based systems (e.g., \expword{early versions of SYSTRAN}).
	\end{itemize}
	
\end{frame}

\begin{frame}
	\frametitle{\insertshortsubtitle: \insertsection}
	\framesubtitle{\insertsubsection: Transfer-based}
	
	\begin{itemize}
		\item \textbf{Algorithm}:
		\begin{itemize}
			\item Compute the parse tree of the source text $S$ (analysis stage).
			\item Apply lexical and structural transfer rules to map the source representation into a target-language representation.
			\item Generate the target text $T$ from the resulting parse tree (generation stage).
		\end{itemize}
		\item \textbf{Preconditions}:
		\begin{itemize}
			\item Grammars of both source and target languages;
			\item Transfer rules (manually defined or learned from corpora);
			\item A high-quality bilingual dictionary;
			\item Tools for morphological analysis and parsing (source language);
			\item Tools for text generation from syntactic structures (target language).
		\end{itemize}
		\item \textbf{Level}: Lexical and syntactic.
	\end{itemize}
	
\end{frame}


\begin{frame}
	\frametitle{\insertshortsubtitle: \insertsection}
	\framesubtitle{\insertsubsection: Transfer-based (Example of transfer rules)}

	\begin{figure}
		\centering
		\hgraphpage[\textwidth]{MT-transfer-exp.pdf}
		\caption{Example of syntactic transfer rules \cite{06-quah}}
	\end{figure}

\end{frame}

\begin{frame}
	\frametitle{\insertshortsubtitle: \insertsection}
	\framesubtitle{\insertsubsection: Transfer-based (Apertium \cite{11-forcada-al})}
	
	\begin{itemize}
		\item Apertium: an open-source rule-based machine translation platform.
		\item Website: \url{https://www.apertium.org/}
	\end{itemize}

	\hgraphpage{MT-apertium-arch.pdf}
	
\end{frame}

\begin{frame}
	\frametitle{\insertshortsubtitle: \insertsection}
	\framesubtitle{\insertsubsection: Interlingua-based}
	
	\begin{figure}
		\centering
		\hgraphpage[.9\textwidth]{MT-Interlingua.pdf}
		\caption{Example of an interlingua-based machine translation model \cite{06-quah}}
	\end{figure}

\end{frame}

\begin{frame}
	\frametitle{\insertshortsubtitle: \insertsection}
	\framesubtitle{\insertsubsection: Interlingua-based (Description)}
	
	\begin{itemize}
		\item \textbf{Algorithm}:
		\begin{itemize}
			\item Analyze the source text $S$ (syntactic and semantic analysis).
			\item Map the analysis into a language-independent representation (interlingua).
			\item Generate a target-language representation from the interlingua.
			\item Produce the target text $T$.
		\end{itemize}
		\item \textbf{Preconditions}:
		\begin{itemize}
			\item Tools for syntactic and semantic analysis;
			\item A well-defined interlingua representation;
			\item Lexical resources linking the interlingua to each language.
		\end{itemize}
		\item \textbf{Level}: Syntactic and semantic.
	\end{itemize}
	
\end{frame}


\begin{frame}[fragile]
	\frametitle{\insertshortsubtitle: \insertsection}
	\framesubtitle{\insertsubsection: Interlingua-based (KANT interlingua \cite{98-czuba-al})}
	
	\begin{columns}
	\begin{column}{0.5\textwidth}
	\bfseries\fontsize{4}{5}\selectfont
\begin{verbatim}
(*A-REMAIN  ; action rep for ’remain’
   (FORM FINITE)
   (TENSE PAST)
   (MOOD DECLARATIVE)
   (PUNCTUATION PERIOD)
   (IMPERSONAL -) ; passive + expletive subject
   (ARGUMENT-CLASS THEME+PREDICATE) ; predicate argument structure
   (Q-MODIFIER ; PP semrole (generic)
      (*K-DURING ; PP interlingua
         (POSITION FINAL) ; clue for translation
            (OBJECT ; PP object semrole
               (*O-TIME ; object rep for ’time’
                  (UNIT -)
                  (NUMBER SINGULAR)
                  (REFERENCE DEFINITE)
                  (DISTANCE NEAR)
                  (PERSON THIRD)))))
   (THEME ; object semrole
      (*O-DEFAULT-RATE ; object rep for ’default rate’
         (PERSON THIRD)
         (UNIT -)
         (NUMBER SINGULAR)
         (REFERENCE DEFINITE)))
   (PREDICATE ; adjective phrase semrole
      (*P-CLOSE ; property rep for ’closer’
         (DEGREE POSITIVE)
         (Q-MODIFIER
            (*K-TO
               (OBJECT
                  (*O-ZERO
                     (UNIT -)
                     (NUMBER SINGULAR)
                     (REFERENCE NO-REFERENCE)
                     (PERSON THIRD))))))))
\end{verbatim}
	\end{column}
	\begin{column}{0.5\textwidth}
		\begin{figure}
			\caption{Sentence ``\textit{The default rate remained close to zero during this time.}'' representation using KANT Interlingua \cite{98-czuba-al}}
		\end{figure}
	\end{column}
	\end{columns}

\end{frame}

\begin{frame}[fragile]
	\frametitle{\insertshortsubtitle: \insertsection}
	\framesubtitle{\insertsubsection: Interlingua-based (KANTOO sys \cite{00-nyberg-al})}
	
	\hgraphpage{MT-kantoo-arch.pdf}
	
\end{frame}

\begin{frame}
	\frametitle{\insertshortsubtitle: \insertsection}
	\framesubtitle{\insertsubsection: Statistical}
	
	\[
	P(T|S) = \frac{P(T) P(S|T)}{P(S)} \propto \underbrace{P(T)}_{\text{fluency}} \underbrace{P(S \mid T)}_{\text{adequacy}}
	\]
	
	\[
	\hat{T} = \arg\max_{T} P(T)\, P(S|T)
	\]
	
	\begin{itemize}
		\item $P(T)$: language model of the target language
		\begin{itemize}
			\item Requires a large corpus in the target language.
			\item $P(T) = \prod_{j=1}^{m} P(t_j \mid t_{j-N+1}, \ldots, t_{j-1})$.
		\end{itemize}
		
		\item $P(S|T)$: translation model
		\begin{itemize}
			\item Requires a parallel (aligned) source--target corpus.
			\item Word alignment problem (different word orders).
			\item Length differences between source and target sentences.
		\end{itemize}
	\end{itemize}
	
\end{frame}

\begin{frame}
	\frametitle{\insertshortsubtitle: \insertsection}
	\framesubtitle{\insertsubsection: Statistical (Alignment)}
	
	\[
	P(S|T) = \sum_{A} P(S, A \mid T)
	\]
	
	\[
	P(S, A \mid T) = \prod_{i=1}^{n} P(S_i, A_i \mid S_1^{i-1}, A_1^{i-1}, T_1^{m})
	\]
	
	\[
	P(S_i, A_i | S_1^{i-1}, A_1^{i-1}, T_1^{m}) = P(A_i | S_1^{i-1}, A_1^{i-1}, T_1^{m}) P(S_i | S_1^{i-1}, A_1^{i}, T_1^{m})
	\]
	
	\begin{itemize}
		\item $A_i$: alignment of source word $S_i$ to a target word (or NULL).
		\item Number of possible alignments: $(m + 1)^n$.
	\end{itemize}
	
\end{frame}

\begin{frame}
	\frametitle{\insertshortsubtitle: \insertsection}
	\framesubtitle{\insertsubsection: Statistical (IBM Model 1)}
	
	\begin{itemize}
		\item IBM Model 1: uniform alignment distribution and independent word generation.
	\end{itemize}
	
	\[
	P(A_i | S_1^{i-1}, A_1^{i-1}, T_1^{m}) = \frac{1}{m+1}
	\]
	
	\[
	P(S|T) = \frac{1}{(m+1)^n} \sum_{A} \prod_{i=1}^{n} P(S_i | S_1^{i-1}, A_1^{i}, T_1^{m})
	\]
	
\end{frame}

\begin{frame}
	\frametitle{\insertshortsubtitle: \insertsection}
	\framesubtitle{\insertsubsection: Statistical (IBM Model 2)}
	
	\begin{itemize}
		\item IBM Model 2 introduces position-based alignment probabilities.
		\item Alignment probability $P(A_i \mid i, n, m)$ depends only on:
		\begin{itemize}
			\item the position $i$ of the source word;
			\item the aligned target word index $A_i$;
			\item the source sentence length $n$ and target sentence length $m$.
		\end{itemize}
		\item This probability is estimated from a parallel corpus using relative frequencies.
	\end{itemize}
	
	\[
	P(A_i | S_1^{i-1}, A_1^{i-1}, T_1^{m}) = P(A_i | i, n, m)
	\]
	
	\[
	P(S|T) = \sum_{A} \prod_{i=1}^{n} P(A_i | i, n, m) P(S_i | S_1^{i-1}, A_1^{i}, T_1^{m})
	\]
	
\end{frame}

\begin{frame}
	\frametitle{\insertshortsubtitle: \insertsection}
	\framesubtitle{\insertsubsection: Statistical (HMM alignment)}
	
	\begin{itemize}
		\item HMM alignment model.
		\item Alignment of a source word $S_i$ depends only on the previous alignment $A_{i-1}$ and the target sentence length $m$.
		\item Lexical generation follows IBM1: $S_i$ depends only on the aligned target word $T_{A_i}$
	\end{itemize}
	
	\[
	P(A_i | S_1^{i-1}, A_1^{i-1}, T_1^{m}) = P(A_i | A_{i-1}, m)
	\]
	
	\[
	P(S|T) = \sum_{A} \prod_{i=1}^{n} P(A_i | A_{i-1}, m) P(S_i | S_1^{i-1}, A_1^{i}, T_1^{m})
	\]
	
\end{frame}

\begin{frame}
	\frametitle{\insertshortsubtitle: \insertsection}
	\framesubtitle{\insertsubsection: Example-based MT}
	
	\begin{itemize}
		\item EBMT relies on a database of bilingual examples rather than purely probabilistic models.
		\item Translation is performed on segments (contiguous sequences of words) instead of individual words.
		\item Segment translation probability is estimated by comparing source segments with stored target segments.
		\item Phrase-based MT (PBMT):
		\begin{itemize}
			\item Uses statistical models trained on bilingual corpora.
			\item Breaks sentences into phrases (short sequences of words).
			\item Computes probabilities for phrase translations and reorderings.
			\item Combines phrases to produce the most probable target sentence.
			\item Optimizes fluency using a target-language model.
		\end{itemize}
	\end{itemize}
	
\end{frame}

\begin{frame}
	\frametitle{\insertshortsubtitle: \insertsection}
	\framesubtitle{\insertsubsection: Phrase-based MT (Moses system \cite{07-koehn-al})}
	
	\begin{itemize}
		\item Official website: \url{http://statmt.org/moses/}
		\item Segment translation table $\phi(S|T)$: stores source segment $S$, target segment $T$, and a translation probability.
		\item Language model $LM$: trained on target language to ensure fluency.
		\item Distortion model $D(T,S)$: assigns a cost for reordering segments.
		\item Word penalty $W(T)$: penalizes translations that are too long or too short.
		\item Beam search is used to efficiently find the most probable translation $\hat{T}$.
	\end{itemize}
	
	\[
	pP(T|S) = \phi(S|T)^{\lambda_\phi} \times LM(T)^{\lambda_{LM}} \times D(T,S)^{\lambda_D} \times W(T)^{\lambda_W}
	\]
	
\end{frame}

\begin{frame}
	\frametitle{\insertshortsubtitle: \insertsection}
	\framesubtitle{\insertsubsection: Neural Machine Translation (NMT)}
	
	\begin{itemize}
		\item Structure
		\begin{itemize}
			\item Model: encoder-decoder (seq2seq)
			\item \textbf{Encoder}: maps the source sentence $S$ to a context vector representing sentence semantics.
			\item \textbf{Decoder}: generates the target sentence $T$ from the context vector.
		\end{itemize}
		
		\item Formal model
		\[
		P(T|S) = P(t_1|S) P(t_2|S, t_1) P(t_3|S, t_1, t_2)\ldots P(t_m|S, t_1\ldots t_{m-1})
		\]
		\[
		\hat{T} = \arg\max_{T} \prod_{i=1}^{m} P(t_i | S, t_1\ldots t_{i-1})
		\]
	\end{itemize}
	
\end{frame}


\begin{frame}
	\frametitle{\insertshortsubtitle: \insertsection}
	\framesubtitle{\insertsubsection: NN-based (Google translate \cite{16-wu-al})}
	
	\begin{itemize}
		\item Google Translate: first large-scale NN-based MT system
		\item Uses encoder-decoder with attention and large bilingual corpora
		\item Accessible at \url{https://translate.google.com/}
	\end{itemize}
	
	\begin{center}
		\hgraphpage[.75\textwidth]{MT-googlet.pdf}
	\end{center}

\end{frame}

\begin{frame}
	\frametitle{\insertshortsubtitle: \insertsection}
	\framesubtitle{\insertsubsection: NN-based (OpenNMT system \cite{17-klein-al})}
	
	\begin{itemize}
		\item OpenNMT: open-source neural MT framework
		\item Supports seq2seq, attention, and Transformer models
		\item Website: \url{https://opennmt.net/}
	\end{itemize}
	
	\begin{center}
		\hgraphpage[.85\textwidth]{MT-opennmt.pdf}
	\end{center}
	
\end{frame}

\begin{frame}
	\frametitle{\insertshortsubtitle: \insertsection}
	\framesubtitle{\insertsubsection: Some humor}
	
	\begin{center}
		\hgraphpage[0.65\textwidth]{humor/humor-translation1.jpg}
		\hgraphpage[0.25\textwidth]{humor/humor-translation2.jpeg}
	\end{center}
	
\end{frame}

\subsection{Automatic text summarization}

\begin{frame}
	\frametitle{\insertshortsubtitle: \insertsection}
	\framesubtitle{\insertsubsection: Summary categories}
	
	\hgraphpage{ATS-classif.pdf}
	
\end{frame}

\begin{frame}
	\frametitle{\insertshortsubtitle: \insertsection}
	\framesubtitle{\insertsubsection: Approaches}
	
	\begin{columns}
		\begin{column}{0.4\textwidth}
			\begin{block}{\scriptsize\bfseries\cite{12-nenkova-mckeown}}
				\begin{itemize}
					\item Topic representation:
					\begin{itemize}
						\item topic words
						\item frequencies
						\item latent semantic analysis
						\item Bayesian topic models
						\item clustering
					\end{itemize}
					\item Indicators representation: 
					\begin{itemize}
						\item graph-based
						\item machine learning (ML)
					\end{itemize}
				\end{itemize}
			\end{block}
		\end{column}
		\begin{column}{0.3\textwidth}
			\begin{block}{\scriptsize\bfseries\cite{12-lloret-palomar}}
				\begin{itemize}
					\item statistical 
					\item graph-based
					\item discourse-based
					\item ML-based
				\end{itemize}
			\end{block}
		\end{column}
		\begin{column}{0.28\textwidth}
			\begin{block}{\scriptsize\bfseries\cite{19-aries-al}}
				\begin{itemize}
					\item statistical
					\item graph-based
					\item linguistic 
					\item ML-based
				\end{itemize}
			\end{block}
		\end{column}
	\end{columns}

\end{frame}

\begin{frame}
	\frametitle{\insertshortsubtitle: \insertsection}
	\framesubtitle{\insertsubsection: Statistical features}
	
	\begin{itemize}
		\item \optword{Word frequency} 
		
		\hspace{.5cm}E.g. $Score_\text{TF-IDF}(s_i) = \sqrt{\sum\limits_{w_{ik} \in s_i} (\text{TF-IDF}(w_{ik}))^2}$
		
		\item \optword{Sentence (or word) Position}
		
		\hspace{.5cm}E.g. $ Score_\text{position}(s_i) = \max (\frac{1}{i}, \frac{1}{|D| - i + 1}) $
		
		\item \optword{Sentence length}
		
		\hspace{.5cm}E.g. $ Score_\text{length}(s_i) = \left\lbrace 
		\begin{array}{lll}
		0 & si & (L_i \geq L_{\text{min}}) \\
		\frac{L_i - L_{\text{min}}}{L_{\text{min}}} & \text{otherwise} & \\
		\end{array}
		\right. $
		
		\item \optword{Title and subtitle words}
		
		\hspace{.5cm}E.g. $ Score_{\text{title}}(s_i) = \frac{\sum_{e \in T \bigcap s_i}{\frac{tf(e)}{tf(e)+1}}}
		{\sum_{e \in T}{\frac{tf(e)}{tf(e)+1}}} $
		
		\item \optword{Centroid}, \optword{Frequent itemsets}, \optword{LSA}, ...
	\end{itemize}

\end{frame}

\begin{frame}
	\frametitle{\insertshortsubtitle: \insertsection}
	\framesubtitle{\insertsubsection: Statistical (e.g., \cite{13-aries-al} (1))}
	
	\hgraphpage{ATS-tcc-arch.pdf}
	
\end{frame}

\begin{frame}
	\frametitle{\insertshortsubtitle: \insertsection}
	\framesubtitle{\insertsubsection: Statistical (e.g., \cite{13-aries-al} (2))}
	
	\begin{itemize}
		\item A text can discuss multiple topics, as well as a single sentence
		\begin{itemize}
			\item Clustering sentences based on similarity with a threshold ($Th$)
		\end{itemize}
		\item A sentence is considered pertinent if it represents the maximum number of clusters
		\[ Score(s_i , c_j , f_k ) = 1 + \sum_{\phi \in s_i} {P(f_k=\phi | s_i \in c_j)} \]
		\[ Score(s_i , \bigcap_{j} c_j , F) =  %\propto 
		\prod_{j} \prod_{k} Score(s_i , c_j , f_k ) \]
		$ s $: sentence, $ c $: cluster, $ f $: feature, $ F $: set of features, $ \phi $: observation of $ f $.
		\item $f$: TF (Uni, Bi); Pos (interval of 10); Len (all words, filtered)
	\end{itemize}
	
\end{frame}

\begin{frame}
	\frametitle{\insertshortsubtitle: \insertsection}
	\framesubtitle{\insertsubsection: Statistical (e.g., \cite{15-oufaida-al} (1))}
	
	\begin{itemize}
		\item Words representation: pre-trained embeddings (Polyglot).
		\item Clustering to extract text's sub-subjects.
		\begin{itemize}
			\item For each word $w_i \in S_1$, find the most similar word $w_j \in S_2$:
			\[
			Match(w_i | S_2) = \arg\max_{w_j \in S_2} sim(Rep(w_i), Rep(w_j))
			\]
			\item Sentence similarity between $S_1$ and $S_2$:
			{\scriptsize
			\[
			Sim(S_1, S_2) = 
			\]
			\[
			\frac{\sum_{w_i \in S_1} sim(Rep(w_i), Rep(Match(w_i|S_2))) + \sum_{w_j \in S_2} sim(Rep(w_j), Rep(Match(w_j|S_1)))}{|S_1| + |S_2|}
			\]
			}
			\item Apply a clustering algorithm on the sentence similarity matrix to identify sub-subjects.
		\end{itemize}
	\end{itemize}
	
\end{frame}

\begin{frame}
	\frametitle{\insertshortsubtitle: \insertsection}
	\framesubtitle{\insertsubsection: Statistical (e.g., \cite{15-oufaida-al} (2))}
	
	\begin{itemize}
		\item Term scoring based on mRMR
		\begin{itemize}
			\item Mutual information (X, Y: term/sentence vectors) 
			\[
			I(X, Y) = \sum_{x \in X} \sum_{y \in Y} p(x, y) \log \frac{p(x, y)}{p(x) p(y)}
			\]
			
			\item Term relevance (H: cluster/sentence vector) $Relevance(T_i) = I(T_i, H)$
%			\[
%			Relevance(T_i) = I(T_i, H)
%			\]
			
			\item Term redundancy (S: all sentences) $Redundancy(T_i) = \frac{1}{|S|} \sum_{j \in S} I(T_i, T_j)$
%			\[
%			Redundancy(T_i) = \frac{1}{|S|} \sum_{j \in S} I(T_i, T_j)
%			\]
			
			\item Final term score: mRMR vector
			\[
			MID \equiv \max_{t \in T} \big(Relevance(t) - Redundancy(t)\big); \,\, MIQ \equiv \max_{t \in T} \frac{Relevance(t)}{Redundancy(t)}
			\]
%			\[
%			MIQ \equiv \max_{t \in T} \frac{Relevance(t)}{Redundancy(t)}
%			\]
		\end{itemize}
		
		\item A sentence is scored based on its similarity to the mRMR vector.  
		\item After selecting a sentence, the mRMR vector is updated by decreasing the weights of the selected terms.
	\end{itemize}
	
\end{frame}


\begin{frame}
	\frametitle{\insertshortsubtitle: \insertsection}
	\framesubtitle{\insertsubsection: Graph-based}
	
	\begin{itemize}
		\item \optword{Graph properties}
		\begin{itemize}
			\item Bushy paths 
			
			\hspace{.5cm}$Score_{\#arcs}(s_i) = |\{ s_j : a(s_i, s_j) \in A / s_j \in S, s_i \neq s_j \}|$
			
			\item Aggregate similarity
			
			\hspace{.5cm}$Score_{aggregate}(s_i) = \sum\limits_{(s_i, s_j) \in E} sim(s_i, s_j)$
		\end{itemize}
		\item \optword{Iterative methods}
		\begin{itemize}
			\item Update each node's score based on its neighbors' scores.
			\item Stopping criterion: convergence to a stable state.
			\item Example: TextRank \cite{04-mihalcea-tarau}
			
			$WS(V_i) = ( 1 - d) + d * \sum\limits_{V_j \in In(V_i)} \frac{w_{ji}}{\sum\limits_{V_k \in Out(V_j)} w_{jk}} WS(V_j)$
			
			$w_{ij} = \frac{|\{w_k \text{ / } w_k \in S_i \text{ and } w_k \in S_j\}|}{\log(|S_i|) + \log(|S_j|)}$
			
			\item \( d \) is the damping factor (commonly \(0.85\)).
		\end{itemize}
	\end{itemize}

\end{frame}


\begin{frame}
	\frametitle{\insertshortsubtitle: \insertsection}
	\framesubtitle{\insertsubsection: Graph-based (e.g., \cite{21-aries-al} (1))}
	
	\begin{center}
		\vgraphpage{ATS-gc-archi.pdf}
	\end{center}
	
\end{frame}

\begin{frame}
	\frametitle{\insertshortsubtitle: \insertsection}
	\framesubtitle{\insertsubsection: Graph-based (e.g., \cite{21-aries-al} (2))}
	
	\vspace{-6pt}
	\begin{itemize}
		\item Graph simplification
		
		\hspace{.5cm}$weak\_node(v_i) = ( \sum_{(v_i, v_j) \in E} w_{ij} < \frac{1}{MImpN(v_i)} )$ 
		
		\hspace{.5cm}$weak\_arc(v_i, v_j) = ( w_{ij} < \frac{Threshold}{MImpN(v_i)})$
		
		\item Statistical score of a sentence
		
		\hspace{.5cm}$ Score(s_i/ sim) = sim(s_i, C\backslash s_i) $
		$Score(s_i/ tfisf) = \sqrt{\sum\limits_{w_{ik} \in s_i} (tfisf(w_{ik}))^2}$
		
		\hspace{.5cm}$Score(s_i/ size) = \frac{1}{|s_i|}$
		$Score(s_i/ pos) = \max (\frac{1}{i}, \frac{1}{|D| - i + 1})$
		
		\hspace{.5cm}$SSF(s_i/ F) = \prod_{f_i \in F} score(s_i/f_i)$
		
		\item Graph-based score of a sentence
		
		\hspace{.5cm}E.g. $GC1(s_i) = SSF(s_i) + \sum\limits_{(s_i, s_j) \in E} sim(s_i, s_j) * SSF(s_j)$
		
		\item Extraction
		
		\hspace{.5cm}E.g. $ next_{e4}  =  \arg\min\limits_i (iord\ gc(s_i) + ord\ sim(last_{e4}, i))$ 
		
		\hspace{3cm}$ \text{ where } (last_{e4}, s_i) \in E $
	\end{itemize}
	
\end{frame}

\begin{frame}
	\frametitle{\insertshortsubtitle: \insertsection}
	\framesubtitle{\insertsubsection: Linguistic}
	
	\begin{itemize}
		\item \optword{Subject words}: a list of subject-pertinent words, such as ``significant", ``impossible", etc.
		\begin{itemize}
			\item E.g. \cite{69-edmundson}: Bonus (positively pertinent words), Stigma (negatively pertinent words).
			
			$Score_{cue}(s_i) = \sum_{w \in s_i}{cue(w)}
			\text{ where }
			cue(w) = \left\lbrace 
			\begin{array}{ll}
			b > 0 & \text{if } (w \in Bonus) \\
			\delta < 0 & \text{if } (w \in Stigma) \\
			0 & \text{otherwise} 
			\end{array} 
			\right. $
		\end{itemize}
		\item \optword{Indicators}: Structures that imply that the sentence containing them has something important about the subject.
		\begin{itemize}
			\item E.g. \expword{the principal aim of this paper is to investigate ...}
		\end{itemize}
		\item \optword{Co-reference}: use of anaphora resolution or semantic resources (e.g., WordNet).
		\item \optword{Rhetorical structure}: use of rhetorical structure to score sentences or phrases.
	\end{itemize}
	
\end{frame}

\begin{frame}
	\frametitle{\insertshortsubtitle: \insertsection}
	\framesubtitle{\insertsubsection: Linguistic (e.g.: \cite{81-paice})}
	
	\begin{itemize}
		\item Using manually prepared templates
		\item\ [x]: up to x words can be between he current word and the next
		\item +y: update the score by y
		\item ?: the current word is optional
	\end{itemize}

	\begin{figure}[!ht]
		\begin{center}
			\hgraphpage[.7\textwidth]{ATS-paice-template.pdf}
			\caption{Example of a simplified template \cite{81-paice}.}
			\label{fig:paice-template}
		\end{center}
	\end{figure}
	
\end{frame}

\begin{frame}
	\frametitle{\insertshortsubtitle: \insertsection}
	\framesubtitle{\insertsubsection: ML-based}
	
	\begin{itemize}
		\item \optword{Features-based}
		\begin{itemize}
			\item Hyperparameter tuning: adjust statistical scores' weights
			\item Classification: decide if a text unit (sentence) belongs to the summary
		\end{itemize}
		
		\item \optword{Bayesian topic models}
		\begin{itemize}
			\item Identify principal concepts from documents and links between them to create a hierarchy
		\end{itemize}
		
		\item \optword{Deep learning}
		\begin{itemize}
			\item Hyperparameter tuning or sentence classification
			\item Word-level generation (can be framed as classification)
		\end{itemize}
		
		\item \optword{Reinforcement learning}
		\begin{itemize}
			\item Use actions and rewards to train a system to generate summaries
		\end{itemize}
	\end{itemize}
	
\end{frame}

\begin{frame}
	\frametitle{\insertshortsubtitle: \insertsection}
	\framesubtitle{\insertsubsection: ML-based (e.g., \cite{2020-aries})}
	
	\begin{center}
		\vgraphpage{ATS-ml2es-archi.pdf}
	\end{center}
	
\end{frame}

\begin{frame}
	\frametitle{\insertshortsubtitle: \insertsection}
	\framesubtitle{\insertsubsection: ML-based (e.g., \cite{06-daumeiii-marcu})}
	
	\begin{minipage}{.6\textwidth}
		\begin{itemize}
			\item $D$: a set of $K$ documents
			\item $Q$: a set of $J$ queries
			\item $P^G$: a language model trained on general English 
			\item $P^Q$: a language model trained on requests
			\item $P^D$: a language model trained on documents
		\end{itemize}
	\end{minipage}
	\begin{minipage}{.38\textwidth}
		\hgraphpage{ATS-btm-daumeiii.pdf}
	\end{minipage}
	
\end{frame}

\begin{frame}
	\frametitle{\insertshortsubtitle: \insertsection}
	\framesubtitle{\insertsubsection: ML-based (e.g., \cite{15-rush-al})}
	
	\begin{center}
		\hgraphpage[.35\textwidth]{ATS-2015-rush-al-arch.pdf}
		\hgraphpage[.35\textwidth]{ATS-2015-rush-al-exp.pdf}
	\end{center}
	
	\begin{itemize}
		\item[(a)] A decoder which generates the next word $y_{i+1}$ of the summary based on the input text $x$ and the words already generated $y_c \equiv [y_{i-c+1},\ldots, y_i]$
		\item[(b)] An encoder with attention
	\end{itemize}
	
\end{frame}

\begin{frame}
	\frametitle{\insertshortsubtitle: \insertsection}
	\framesubtitle{\insertsubsection: ML-based (e.g., \cite{18-narayan-al})}
	
	\hgraphpage{ATS-narayan-al.pdf}
	
\end{frame}

\begin{frame}
	\frametitle{\insertshortsubtitle: \insertsection}
	\framesubtitle{\insertsubsection: Some humor}
	
	\begin{center}
		\hgraphpage[0.6\textwidth]{humor/humor-summarize1.jpg}
		\hgraphpage[0.3\textwidth]{humor/humor-summarize2.jpg}
	\end{center}
	
\end{frame}

%===================================================================================
\section{Interaction}
%===================================================================================

\begin{frame}
	\frametitle{\insertshortsubtitle}
	\framesubtitle{\insertsection}
	
	\begin{itemize}
		\item Tasks where a system takes a user request and interacts to generate an appropriate textual response.
		\item \textbf{Input}: user request
		\item \textbf{Output}: system response
		\item \textbf{Applications:}
		\begin{itemize}
			\item Question Answering
			\item Dialogue systems
		\end{itemize}
	\end{itemize}

\end{frame}

\subsection{Question-Answering}

\begin{frame}
	\frametitle{\insertshortsubtitle: \insertsection}
	\framesubtitle{\insertsubsection}

	\hgraphpage{QA-classif.pdf}
\end{frame}

\begin{frame}
	\frametitle{\insertshortsubtitle: \insertsection}
	\framesubtitle{\insertsubsection: IR-based}
	
	\begin{figure}
		\hgraphpage{QA-IR.pdf}
		\caption{Architecture of a question/answering system \cite{2019-jurafsky-martin}}
	\end{figure}
	
\end{frame}

\begin{frame}
	\frametitle{\insertshortsubtitle: \insertsection}
	\framesubtitle{\insertsubsection: IR-based (Question processing)}
	
	\begin{itemize}
		\item Request formulation
		\begin{itemize}
			\item Word tokenization
			\item Stop words filtering
			\item Stemming
		\end{itemize}
		\item Answer type detection
		\begin{itemize}
			\item Identify entity type: person, place, organization, etc.
			\item Using taxonomies such as WordNet for semantic classification.
		\end{itemize}
	\end{itemize}
	
\end{frame}

\begin{frame}
	\frametitle{\insertshortsubtitle: \insertsection}
	\framesubtitle{\insertsubsection: IR-based (Search)}
	
	\begin{itemize}
		\item Document retrieval:
		\begin{itemize}
			\item Using keywords and document indices;
			\item Return the highest-scoring documents;
			\item Split these documents into passages (paragraphs or sentences).
		\end{itemize}
		\item Passage retrieval:
		\begin{itemize}
			\item Apply NER on passages to identify candidate entities matching the expected answer type;
			\item Filter passages that do not contain the expected answer type;
			\item ML algorithms can score and rank passages based on relevance;
			\item Example features for passage scoring:
			\begin{itemize}
				\item Number of named entities matching the answer type,
				\item Number of terms in the question present in the passage,
				\item Length of the longest sequence similar to the question,
				\item Passage position within the document.
			\end{itemize}
		\end{itemize}
	\end{itemize}
	
\end{frame}

\begin{frame}
	\frametitle{\insertshortsubtitle: \insertsection}
	\framesubtitle{\insertsubsection: IR-based (Answer extraction)}
	
	\begin{itemize}
		\item \optword{Features-based ML}
		\begin{itemize}
			\item Patterns: e.g., \expword{\textless AP\textgreater such as \textless QP\textgreater};
			\item Matched answer type;
			\item Number of question keywords matched;
			\item Novelty factor: presence of at least one word not in the question;
			\item Punctuation;
			\item Other syntactic or lexical features.
		\end{itemize}
		\item \optword{Neural Networks}
		\begin{itemize}
			\item Model performs a reading comprehension task;
			\item For each word, compute the probability of being the start or end of the answer span, conditioned on the question.
		\end{itemize}
	\end{itemize}
	
\end{frame}

\begin{frame}
	\frametitle{\insertshortsubtitle: \insertsection}
	\framesubtitle{\insertsubsection: IR-based (Extraction using Bi-LSTM \cite{2017-chen-al})}
	
	\vskip-6pt
	\begin{figure}
		\vgraphpage[0.8\textheight]{QA-bilstm-exp.pdf}\vskip-3pt
		\caption{Answer extraction in the DrQA system \cite{2019-jurafsky-martin}}
	\end{figure}
	
\end{frame}

\begin{frame}
	\frametitle{\insertshortsubtitle: \insertsection}
	\framesubtitle{\insertsubsection: IR-based (Extraction using BERT \cite{2018-devlin-al})}
	
	\begin{figure}
		\centering
		\hgraphpage[.5\textwidth]{QA-bert-exp.pdf}
		\caption{Answer span extraction using BERT \cite{2019-jurafsky-martin}}
	\end{figure}
	
\end{frame}

\begin{frame}
	\frametitle{\insertshortsubtitle: \insertsection}
	\framesubtitle{\insertsubsection: Knowledge-based}
	
	\begin{itemize}
		\item \optword{Graph-based}:
		\begin{itemize}
			\item Apply entity linking to identify entities in the question;
			\item Determine the target relation (e.g., \expword{Birth\_place});
			\item If multiple candidate relations exist: calculate similarity between the question and candidate relations;
			\item Similarity can be computed using embeddings.
		\end{itemize}
		\item \optword{Semantic-parsing-based}:
		\begin{itemize}
			\item Parse the question into a structured representation;
			\item Representation can be in lambda calculus, SQL, SPARQL, etc.;
			\item This structured form is used as a query to a knowledge base;
			\item Retrieve the answer from the knowledge base.
		\end{itemize}
	\end{itemize}
	
\end{frame}


\begin{frame}
	\frametitle{\insertshortsubtitle: \insertsection}
	\framesubtitle{\insertsubsection: Knowledge-based (Example)}
	
	\begin{exampleblock}{Example of logical forms \cite{2020-jurafsky-martin}}
		\centering\tiny\bfseries
		\begin{tabular}{p{0.5\textwidth}p{0.45\textwidth}}
			\hline\hline
			Question & Logical form \\
			\hline
			Which Wilayas border Jijel?  
			& $\lambda x.wilaya(x) \wedge border(x, Jijel)$ \\
			
			What is the largest Wilaya? 
			& $\arg\max(\lambda x.wilaya(x), \lambda x.area(x))$ \\
			
			Book a flight from Jijel to Algiers.
			& {\color{blue}SELECT DISTINCT} f1.flight id
			
			{\color{blue}FROM} vol v1, airport  a1,
			
			wilaya w1, airport\_service a2, wilaya w2
			
			{\color{blue}WHERE} v1.depart=a1.airport\_code
			
			{\color{blue}AND} a1.wilaya\_code=w1.wilaya\_code
			
			{\color{blue}AND} w1.wilaya\_name= '{\color{red}Jijel}'
			
			{\color{blue}AND} v1.destination=a2.airport\_code
			
			{\color{blue}AND} a2.wilaya\_code=w2.wilaya\_code
			
			{\color{blue}AND} w2.wilaya\_name= '{\color{red}Algiers}' \\
			
			
			How many people survived the sinking of Titanic ? 
			& (count (!fb:event.disaster.survivors
			
			fb:en.sinking of the titanic))\\
			
			How many yards longer was Johnson's longest touchdown compared to his shortest touchdown of the first quarter?
			& 
			ARITHMETIC diff( SELECT num( ARGMAX(
			
			SELECT ) ) SELECT num( ARGMIN( FILTER(
			
			SELECT ) ) ) )\\
			\hline\hline
		\end{tabular}
	\end{exampleblock}
	
\end{frame}

\begin{frame}
	\frametitle{\insertshortsubtitle: \insertsection}
	\framesubtitle{\insertsubsection: LM-based}
	
	\begin{itemize}
		\item Employ a pretrained language model.
		\item Fine-tune the model for a question-answering task.
	\end{itemize}
	
	\begin{figure}
		\centering
		\hgraphpage[.6\textwidth]{qa-t5.pdf}
		\caption{T5 is pretrained on a text infilling task and subsequently fine-tuned for question-answering, without requiring additional context beyond the input question. \cite{2020-roberts-al}}
	\end{figure}
	
\end{frame}


\begin{frame}
	\frametitle{\insertshortsubtitle: \insertsection}
	\framesubtitle{\insertsubsection: Some humor}
	
	\begin{columns}
		\begin{column}{0.3\textwidth}
			\hgraphpage{humor/humor-QR1.jpg}
			
			\hgraphpage{humor/humor-QR3.jpeg}
		\end{column}
		\begin{column}{0.4\textwidth}
			\hgraphpage{humor/humor-QR2.jpeg}
		\end{column}
	\end{columns}

\end{frame}

\subsection{Dialogue systems}

\begin{frame}
	\frametitle{\insertshortsubtitle: \insertsection}
	\framesubtitle{\insertsubsection}
	
	\hgraphpage{DS-classif.pdf}
	
\end{frame}

\begin{frame}
	\frametitle{\insertshortsubtitle: \insertsection}
	\framesubtitle{\insertsubsection: Task-oriented - Frame-based (Frame)}
	
	\begin{itemize}
		\item Frame: a structure containing slots to be filled and predefined questions for each slot.
		\item Only questions corresponding to unfilled slots are asked.
		\item Slots in other frames can also be filled automatically. 
		E.g., \expword{The slot RESERVATION\_DATE of the frame HOTEL\_RESERVATION can be filled from the slot ARRIVAL\_DATE of the frame FLIGHT\_RESERVATION}.
		\item Use \keyword{Intention detection} to identify the relevant frame for the user request.
	\end{itemize}
	
	\vspace{-6pt}
	\begin{exampleblock}{Example of flight frame \cite{2020-jurafsky-martin}}
		\centering\tiny\bfseries
		\begin{tabular}{lll}
			\hline\hline
			Slot & Type & Question \\
			\hline
			ORIGIN CITY & city & ``From what city are you leaving?" \\
			DESTINATION CITY & city & ``Where are you going?" \\
			DEPARTURE TIME & time & ``When would you like to leave?" \\
			DEPARTURE DATE & date & ``What day would you like to leave?" \\
			ARRIVAL TIME & time & ``When do you want to arrive?" \\
			ARRIVAL DATE & date & ``What day would you like to arrive?" \\
			\hline\hline
		\end{tabular}
	\end{exampleblock}
	
\end{frame}

\begin{frame}
	\frametitle{\insertshortsubtitle: \insertsection}
	\framesubtitle{\insertsubsection: Task-oriented - Frame-based (Multiple answers)}
	
		\centering\footnotesize
		\begin{tabular}{lll}
			SHOW & \textrightarrow & show me \textbar\ I want \textbar\ can-I see \\
			DEPARTTIME & \textrightarrow & (after \textbar\ around \textbar\ before) HOUR \textbar\\
			&  & morning \textbar\ afternoon \textbar\ evening\\
			DEPARTDATE & \textrightarrow & on (Saturday \textbar\ ... \textbar\ Friday)\\
			HOUR & \textrightarrow & (one \textbar\ two \textbar\ ... \textbar\ twelve) (am \textbar\ pm) \\
			FLIGHTS & \textrightarrow & (a) flight \textbar\ flights \\
			ORIGIN & \textrightarrow & from WILAYA \\
			DESTINATION & \textrightarrow & to WILAYA \\
			WILAYA & \textrightarrow & Adrar \textbar\ ... \textbar\ El Meniaa \\
		\end{tabular}
		
		\hgraphpage[.8\textwidth]{DS-frame-parse-exp.pdf}
	
\end{frame}

\begin{frame}
	\frametitle{\insertshortsubtitle: \insertsection}
	\framesubtitle{\insertsubsection: Task-oriented - Dialogue-State}
	
	\begin{figure}
		\centering
		\hgraphpage{DS-dialog-arch.pdf}
		\caption{Architecture of a system using dialogue-state \cite{2016-williams-al}}
	\end{figure}
	
\end{frame}

\begin{frame}
	\frametitle{\insertshortsubtitle: \insertsection}
	\framesubtitle{\insertsubsection: Task-oriented - Dialogue-State (Dialogue acts)}
	
	\vspace{-4pt}
	\begin{figure}
		\centering\tiny\bfseries
		\begin{tabular}{llll}
			\hline
			Tag & Sys & User & Description \\
			\hline
			HELLO (a = x, b = y, ...) & \CheckedBox & \CheckedBox & Open a dialogue and give info a = x, b = y, ... \\
			INFORM(a = x, b = y, ...) & \CheckedBox & \CheckedBox & Give info a = x, b = y, ... \\
			REQUEST(a, b = x, ...) & \CheckedBox & \CheckedBox & Request value for a given b = x, ... \\
			REQALTS(a = x, ...) & \XBox & \CheckedBox & Request alternative with a = x, ... \\
			CONFIRM(a = x, b = y, ...) & \CheckedBox & \CheckedBox & Explicitly confirm a = x, b = y, ... \\
			CONFREQ (a = x, ..., d) & \CheckedBox & \XBox & Implicitly confirm a = x, ... and request value of d \\
			SELECT(a = x, a = y) & \CheckedBox & \XBox & Implicitly confirm a = x, ... and request value of d \\
			AFFIRM(a = x, b = y, ...) & \CheckedBox & \CheckedBox & Affirm and give further info a = x, b = y, ... \\
			NEGATE(a = x) & \XBox & \CheckedBox & Negate and give corrected a = x \\
			DENY(a = x) & \XBox & \CheckedBox & Deny that a = x \\
			BYE () & \CheckedBox & \CheckedBox & Close a dialogue \\
			\hline
		\end{tabular}
	
		\begin{tabular}{p{0.5\textwidth}p{0.4\textwidth}}
			\hline
			Utterance & Dialogue act\\
			\hline
			U: Hi, I am looking for somewhere to eat.  & hello(task=find, type=restaurant) \\
			S: You are looking for a restaurant. What type of food do you like?  & confreq(type=restaurant, food) \\
			U: I'd like an Italian somewhere near the museum. & inform(food=Italian, near=museum)\\
			S: ``Roma" is a nice Italian restaurant near the museum.  & inform(name=``Roma", type=restaurant, food=Italian, near=museum) \\
			U: Is it reasonably priced? & confirm(pricerange=moderate) \\
			S: Yes, ``Roma" is in the moderate price range.  & affirm(name=``Roma", pricerange=moderate) \\
			U: What is the phone number? & request(phone) \\
			S: The number of ``Roma" is 385456. & inform(name=``Roma", phone=``385456") \\
			U: Ok, thank you goodbye. & bye() \\
			\hline
		\end{tabular}
		
		\caption{Restaurant recommendation using HIS \cite{2010-young-al}}
	\end{figure}
	
\end{frame}

\begin{frame}
	\frametitle{\insertshortsubtitle: \insertsection}
	\framesubtitle{\insertsubsection: Task-oriented - Dialogue-State (Slots filling)}
	
	\begin{itemize}
		\item Classify each sentence according to its intent, domain, and slot.
		\item Extract relevant information to populate slots, e.g., \expword{Using sequences labeling}.
	\end{itemize}
	
	\begin{figure}
		\centering
		\hgraphpage[.4\textwidth]{DS-dialog-fill-exp.pdf}
		\caption{Architecture for slot filling using BERT \cite{2020-jurafsky-martin}}
	\end{figure}
	
\end{frame}

\begin{frame}
	\frametitle{\insertshortsubtitle: \insertsection}
	\framesubtitle{\insertsubsection: Task-oriented - Dialogue-State (Other components)}
	
	\begin{itemize}
		\item Dialogue State Tracker
		\begin{itemize}
			\item Determine the current state of the dialogue frame, i.e., the values filled for each slot.
			\item Predict the user's most recent dialogue act based on their input.
		\end{itemize}
		
		\item Dialogue Policy
		\begin{itemize}
			\item Determine the next system action $A_i$:
			\[\hat{A}_i = \arg\max_{A_i \in A} P(A_i | Frame_{i-1}, A_{i-1}, U_{i-1})\]
		\end{itemize}
		
		\item Natural Language Generation (NLG)
		\begin{itemize}
			\item Generate a textual response from the dialogue act.
			\item An encoder--decoder model (e.g., Seq2Seq or Transformer) can be trained for this task.
		\end{itemize}
	\end{itemize}
	
\end{frame}

\begin{frame}
	\frametitle{\insertshortsubtitle: \insertsection}
	\framesubtitle{\insertsubsection: Chatbot - Rule-based (e.g., ELIZA \cite{1966-Weizenbaum})}
	
	\begin{itemize}
		\item The best-known rule-based chatbot is \keyword{ELIZA}, designed for psychological counseling.
		\item ELIZA uses rules in the form of \textit{pattern--transform tuples}:
		\expword{\small [(.*) YOU (.*) ME]\textsubscript{[Pattern]} \textrightarrow\ [WHAT MAKES YOU THINK I \$2 YOU?]\textsubscript{[Transform]}}
		
		\expword{Example: ``You hate me'' \textrightarrow\ ``WHAT MAKES YOU THINK I HATE YOU?''}
		
		\item Each template is associated with a set of keywords. 
		The most relevant keyword in the input sentence triggers candidate templates.
		
		\item The template whose pattern best matches the input sentence is selected to generate the response.
	\end{itemize}
	
\end{frame}

\begin{frame}
	\frametitle{\insertshortsubtitle: \insertsection}
	\framesubtitle{\insertsubsection: Chatbot - IR-based}
	
	\begin{itemize}
		\item Let $C$ be a corpus of conversations.
		\item The user's input text is treated as a query $q$.
		\item Traditional IR-based response selection:
		\[\text{Response}(q, C) = \arg\max_{r \in C} \frac{q \cdot r}{|q| |r|}\]
		where $q$ and $r$ are represented as TF-IDF vectors.
		\item Neural embedding-based response selection:
		\[
		h_q = BERT_Q(q)[CLS],\quad h_r = BERT_R(r)[CLS]
		\]
		\[
		\text{Response}(q, C) = \arg\max_{r \in C} h_q \cdot h_r
		\]
	\end{itemize}
	
\end{frame}

\begin{frame}
	\frametitle{\insertshortsubtitle: \insertsection}
	\framesubtitle{\insertsubsection: Chatbot - NLG-based}
	
	\begin{itemize}
		\item Use an encoder/decoder model to generate responses:
		\[
		\hat{r}_i = \arg\max_{w \in V} P(w \mid q, r_1, \ldots, r_{i-1})
		\]
		where $q$ is the user query and $r_1, \ldots, r_{i-1}$ are previously generated tokens.
		
		\item The encoder can process a longer context, e.g., previous dialogue turns.
		
		\item Pretrained language models, such as GPT, can be fine-tuned on conversational data.
	\end{itemize}
	
	\vspace{-6pt}
	\begin{figure}
		\centering
		\hgraphpage[.7\textwidth]{DS-chatbot-encdec-exp.pdf}
		\caption{Example of a chatbot based on an encoder/decoder architecture \cite{2020-jurafsky-martin}}
	\end{figure}
	
\end{frame}


\begin{frame}
	\frametitle{\insertshortsubtitle: \insertsection}
	\framesubtitle{\insertsubsection: Chatbot - NLG-based (Examples)}
	
	\begin{itemize}
		\item ChatGPT is a sibling project of InstructGPT \cite{2022-ouyang-al}
		\begin{itemize}
			\item Based on \keyword{GPT-3.5}, a decoder-only Transformer pretrained on next-word prediction, then fine-tuned on conversations provided by human AI trainers playing both roles: user and AI assistant.
			\item Given the user's request, the model predicts the next tokens as the system's response.
			\item Reinforcement Learning from Human Feedback (RLHF) is applied, where the reward is derived from manually scored conversations.
			\item Optimization is performed using Proximal Policy Optimization (PPO) \cite{2017-schulman-al}.
		\end{itemize}
		
		\item Similar open source projects
		\begin{itemize}
			\item \url{https://github.com/lucidrains/PaLM-rlhf-pytorch}
			\item \url{https://github.com/CarperAI/trlx}
		\end{itemize}
	\end{itemize}
	
\end{frame}


\begin{frame}
	\frametitle{\insertshortsubtitle: \insertsection}
	\framesubtitle{\insertsubsection: Some humor}
	
	\begin{center}
		\hgraphpage[0.5\textwidth]{humor/humor-chatbots1.jpeg}
		\hgraphpage[0.4\textwidth]{humor/humor-chatbots2.jpeg}
	\end{center}
\end{frame}


%===================================================================================
\section{Classification}
%===================================================================================

\begin{frame}
	\frametitle{\insertshortsubtitle}
	\framesubtitle{\insertsection}
	
	\begin{itemize}
		\item \textbf{Input}: textual data
		\item \textbf{Output}: a category or label
		\item \textbf{Applications:} 
		\begin{itemize}
			\item Spam detection
			\item Language identification
			\item Readability classification
			\item Sentiment analysis
			\item Humor classification
			\item Author attribution
			\item ...
		\end{itemize}
	\end{itemize}
	
\end{frame}


\subsection{Sentiment analysis}

\begin{frame}
	\frametitle{\insertshortsubtitle: \insertsection}
	\framesubtitle{\insertsubsection}
	
	\hgraphpage{SA-classif.pdf}
	
\end{frame}

\begin{frame}
	\frametitle{\insertshortsubtitle: \insertsection}
	\framesubtitle{\insertsubsection: Knowledge-based}
	
	\begin{itemize}
		\item Identify sentiment-indicator words:
		\begin{itemize}
			\item Adjectives, adverbs, and sometimes verbs or nouns can indicate sentiment.
			\item \optword{Lexicon-based}: using dictionaries such as WordNet to augment the list of indicators (synonyms, antonyms, etc.).
			\item \optword{Corpus-driven}: using co-occurrence with known sentiment words to expand the list.
		\end{itemize}
		\item Assign scores or labels to these words;
		\item Aggregate word scores to determine overall sentiment:
		\begin{itemize}
			\item \textbf{Granularity}: sentence, paragraph, or document level.
			\item Negations can invert sentiment and must be handled explicitly.
			\item Leverage discourse structures, e.g., \keyword{Rhetorical Structure Theory (RST)}.
			\item Example: assign -1 to negative words, +1 to positive words, then sum word polarities to get sentence polarity.
		\end{itemize}
	\end{itemize}
	
\end{frame}

\begin{frame}
	\frametitle{\insertshortsubtitle: \insertsection}
	\framesubtitle{\insertsubsection: ML-based}
	
	\begin{itemize}
		\item \optword{Features-based}:
		\begin{itemize}
			\item Presence and frequency of words;
			\item Part-of-Speech tags;
			\item Opinion words and phrases;
			\item Negation handling.
		\end{itemize}
		\item \optword{Embeddings-based}:
		\begin{itemize}
			\item Use word embeddings to predict the sentiment class.
			\item Fine-tuning BERT (transfer learning):
			\begin{itemize}
				\item Use the ``[CLS]'' token to estimate classes;
				\item ``[CLS]'' token combined with MLP classifier;
				\item Aggregating all hidden states with architectures such as CNN, RNN, or attention-based layers.
			\end{itemize}
		\end{itemize}
	\end{itemize}
	
\end{frame}

\begin{frame}
	\frametitle{\insertshortsubtitle: \insertsection}
	\framesubtitle{\insertsubsection: Hybrid (e.g., \cite{18-bettiche-al} (1))}
	
	\begin{figure}
		\centering
		\hgraphpage[.6\textwidth]{SA-bettiche-al.pdf}
		\caption{Hybrid architecture proposed by \cite{18-bettiche-al} for polarity detection in the Algerian dialect.}
	\end{figure}
	
\end{frame}

\begin{frame}
	\frametitle{\insertshortsubtitle: \insertsection}
	\framesubtitle{\insertsubsection: Hybrid (e.g., \cite{18-bettiche-al} (2)) }
	
	\begin{itemize}
		\item Vocabulary enrichment: 
		\begin{itemize}
			\item Similar words detection (by selecting a representative)
			
			$Ratio = 1 - Levenstein(w_1, w_2)/(|w_1|+|w_2|)$
			
			E.g. \expword{Ratio(kolach, kollach) = 92\%, Ratio(kolach, khlasse) = 38\%.}
			
			\item Calculate semantic orientation of a word $w$
			
			$SO(w) = \sum_{w_p \in V_p} PMI(w, w_p) - \sum_{w_n \in V_n} PMI(w, w_n)$
			
			$PMI (w_1, w_2) = \log \frac{p(w_1, w_2)}{p(w_1)*p(w_2)}$ (estimated from corpus)
		\end{itemize}
		\item Polarity representation:
		\begin{itemize}
			\item Discrete: $polarity(w) = +1 \text{ if } SO(w) > 0, -1 \text{ if } SO(w) < 0$
			\item Continuous: $polarity(w) = SO(w)$
		\end{itemize}
		\item Rule-based version: $polarity(msg) = \text{sign} \sum_{w \in msg} polarity(w)$
		\item ML-based version
		\begin{itemize}
			\item Features: words' polarities
			\item Input word selection using PMI similarity
			\item Many ML algorithms are tested: NB, DT, SVM, etc.
		\end{itemize}
	\end{itemize}
	
	
\end{frame}

\begin{frame}
	\frametitle{\insertshortsubtitle: \insertsection}
	\framesubtitle{\insertsubsection: hybrid (e.g., \cite{18-guellil-al})}
	
	\begin{figure}
		\centering
		\hgraphpage[.45\textwidth]{SA-guellil-al.pdf}
		\caption{Hybrid architecture proposed by \cite{18-guellil-al} to analyze messages in Arabizi.}
	\end{figure}
	
\end{frame}

\begin{frame}
	\frametitle{\insertshortsubtitle: \insertsection}
	\framesubtitle{\insertsubsection: Some humor}
	
	\begin{center}
		\vgraphpage{humor/humor-sentiment.jpg}
	\end{center}
	
\end{frame}

\subsection{Readability}

\begin{frame}
	\frametitle{\insertshortsubtitle: \insertsection}
	\framesubtitle{\insertsubsection}
	
	\hgraphpage{Readability-classif.pdf}
	
\end{frame}

\begin{frame}
	\frametitle{\insertshortsubtitle: \insertsection}
	\framesubtitle{\insertsubsection: Formula (Flesch-Kincaid Grade Level)}
	
	\[
	206.835 - 1.015 (\frac{\text{\slshape total words}}{\text{\slshape total sentences}})
	- 84.6 (\frac{\text{\slshape total syllables}}{\text{\slshape total words}})
	\]
	
	\begin{center}
			\footnotesize
	    \begin{tblr}{
	    		colspec = {p{.15\textwidth}lp{.5\textwidth}},
	    		row{odd} = {lightblue},
	    		row{even} = {lightyellow},
	    		row{1} = {darkblue},
	    	} 
			\bfseries\textcolor{white}{Score} && \bfseries\textcolor{white}{Difficulty}\\
			90-100 && Very easy to read (5th grade). \\
			80-90 && Easy to read (6th grade).\\
			70-80 && Fairly easy to read (7th grade).\\
			60-70 && Plain English (8th \& 9th grade). \\
			50-60 && Fairly difficult to read (10th to 12th grade). \\
			30-50 && Difficult to read (college). \\
			10-30 && Very difficult to read (college graduate). \\
			0-10 && Extremely difficult to read (professional). \\
		\end{tblr}
	\end{center}
	
\end{frame}

\begin{frame}
	\frametitle{\insertshortsubtitle: \insertsection}
	\framesubtitle{\insertsubsection: Formula (Dale--Chall readability formula)}
	
	\[
	0.1579 (\frac{\text{\slshape difficult words}}{\text{\slshape total words}})
	+ 0.0496 (\frac{\text{\slshape total words}}{\text{\slshape total sentences}})
	\]
	
	\begin{center}
		\footnotesize
		\begin{tblr}{
				colspec = {p{.15\textwidth}lp{.5\textwidth}},
				row{odd} = {lightblue},
				row{even} = {lightyellow},
				row{1} = {darkblue},
			} 
			\bfseries\textcolor{white}{Score} && \bfseries\textcolor{white}{Difficulty}\\
			$\le 4.9$ && Grade 4 and below. \\
			5-5.9 && Grades 5--6. \\
			6-6.9 && Grades 7--8.\\
			7-7.9 && Grades 9--10. \\
			8-8.9 && Grades 11--12. \\
			9-9.9 && Grades 13--15 (college). \\
			10+ && Grades 16 and above. \\
		\end{tblr}
	\end{center}
	
\end{frame}

\begin{frame}
	\frametitle{\insertshortsubtitle: \insertsection}
	\framesubtitle{\insertsubsection: Formula (OSMAN \cite{2016-elhaj-rayson})}
	
	\begin{align*}
		Osman = & 200.791 - 1.015 \times (\frac{\text{\slshape total words}}{\text{\slshape total sentences}}) - 24.181 \times \\
		& \\
		&  (\frac{\text{\slshape difficult words + syllables + complex words + Faseeh words}}{\text{\slshape total words}})
	\end{align*}

	\vfill
	
	\begin{itemize}
		\item Difficult words: more than 5 characters without diacritics
		\item Complex words: more than 4 syllables
		\item Faseeh words: complex words with letters \<|', y', w', _d, .z> or ending with \<wn, wA>
	\end{itemize}

\end{frame}

\begin{frame}
	\frametitle{\insertshortsubtitle: \insertsection}
	\framesubtitle{\insertsubsection: ML-based}
	
	\begin{figure}
		\centering
		\hgraphpage{Readability-ML.pdf}
		\caption{Pipeline of reading difficulty estimation using ML \cite{2014-collins}}
	\end{figure}
	
\end{frame}

\begin{frame}
	\frametitle{\insertshortsubtitle: \insertsection}
	\framesubtitle{\insertsubsection: Some humor}
	
	\begin{center}
		\vgraphpage{humor/humor-readability1.jpg}
	\end{center}
	
\end{frame}

%===================================================================================
\section{Speech}
%===================================================================================

\begin{frame}
	\frametitle{\insertshortsubtitle}
	\framesubtitle{\insertsection}
	
	\begin{itemize}
		\item Speech recognition
		\begin{itemize}
			\item \textbf{Input}: speech
			\item \textbf{Output}: text
			\item \textbf{More specific application}: speech to text (STT)
		\end{itemize}
		\item Speech synthesis
		\begin{itemize}
			\item \textbf{Input}: text
			\item \textbf{Output}: speech
			\item \textbf{More specific application}: text to speech (TTS)
		\end{itemize}
	\end{itemize}

\end{frame}

\subsection{Speech recognition}

\begin{frame}
	\frametitle{\insertshortsubtitle: \insertsection}
	\framesubtitle{\insertsubsection: Classification}
	
	\hgraphpage{ASR-classif.pdf}
	
\end{frame}

\begin{frame}
	\frametitle{\insertshortsubtitle: \insertsection}
	\framesubtitle{\insertsubsection: Features extraction (Sampling and Quantization)}
	
	\begin{itemize}
		\item Convert analog signals into digital form.
		\item The signal is divided into short-time frames; each frame is represented as a vector containing the sampled values.
		\item \optword{Sampling}:
		\begin{itemize}
			\item Sampling rate must be at least twice the maximum frequency of the signal (\keyword{Nyquist theorem}).
			\item Typical human speech frequencies are up to 8--10 kHz\\ (sampling rate $\approx$ 20 kHz).
			\item In telephony, the speech bandwidth is limited to 4 kHz\\ (sampling rate $\approx$ 8 kHz).
		\end{itemize}
		\item \optword{Quantization}:
		\begin{itemize}
			\item Amplitudes of samples are stored as integers.
			\item 8-bit representation: -128 to 127; 16-bit: -32768 to 32767.
			\item The value of a sample at time index $n$ is denoted as $x[n]$.
		\end{itemize}
	\end{itemize}
	
\end{frame}

\begin{frame}
	\frametitle{\insertshortsubtitle: \insertsection}
	\framesubtitle{\insertsubsection: Features extraction (Windowing (1))}
	
	\begin{itemize}
		\item Divide the speech signal into short segments using a window.
		\item Each segment is called a \keyword{Frame}.
		\item Each frame is characterized by two parameters: window size and frame stride (shift/offset).
		\item Typical values: window size 20--30 ms, stride 10 ms.
	\end{itemize}
	
	\begin{figure}
		\centering
		\hgraphpage[.5\textwidth]{ASR-windowing-exp.pdf}
		\caption{Example of windowing \cite{2020-jurafsky-martin}}
	\end{figure}
	
\end{frame}

\begin{frame}
	\frametitle{\insertshortsubtitle: \insertsection}
	\framesubtitle{\insertsubsection: Features extraction (Windowing (2))}
	
	\begin{itemize}
		\item To extract a frame $y[n]$, the input signal $s[n]$ is multiplied by a window function $w[n]$:
		\[
		y[n] = w[n]\,s[n]
		\]
		\item The simplest window is the rectangular window:
		\[
		w[n] = \begin{cases}
			1 & \text{if } 0 \le n \le L-1 \\
			0 & \text{otherwise}
		\end{cases}
		\]
		\item This introduces discontinuities at the frame boundaries, leading to spectral leakage in Fourier analysis.
		\item A common solution is to use the Hamming window:
		\[
		w[n] = \begin{cases}
			0.54 - 0.46 \cos \left(\frac{2\pi n}{L-1}\right) & \text{if } 0 \le n \le L-1 \\
			0 & \text{otherwise}
		\end{cases}
		\]
	\end{itemize}
	
\end{frame}


\begin{frame}
	\frametitle{\insertshortsubtitle: \insertsection}
	\framesubtitle{\insertsubsection: Features extraction (Windowing: Example)}
	
	\begin{figure}
		\centering
		\hgraphpage[.5\textwidth]{ASR-windowing2-exp.pdf}
		\caption{Example of Rectangular and Hamming windows \cite{2020-jurafsky-martin}}
	\end{figure}
	
\end{frame}

\begin{frame}
	\frametitle{\insertshortsubtitle: \insertsection}
	\framesubtitle{\insertsubsection: Feature extraction (DFT)}
	
	\begin{itemize}
		\item \optword{Discrete Fourier Transform (DFT)}
		\item Represents the frequency components of a signal (amplitude and phase), i.e., how its energy is distributed across frequencies.
		\item For a frame $x[n]$ of length $N$:
		\[
		X[k] = \sum_{n=0}^{N-1} x[n] e^{-j\frac{2\pi}{N}kn}
		\]
	\end{itemize}
	\vspace{-6pt}
	\begin{figure}
		\centering
		\hgraphpage[.7\textwidth]{ASR-DFT-exp.pdf}
		\caption{(a) A signal portion of the vowel [iy]; (b) its DFT \cite{2020-jurafsky-martin}}
	\end{figure}
	
\end{frame}

\begin{frame}
	\frametitle{\insertshortsubtitle: \insertsection}
	\framesubtitle{\insertsubsection: Feature extraction (Mel filter bank)}
	
	\begin{itemize}
		\item \optword{Mel filter bank}
		\item Human hearing is not equally sensitive across frequency bands.
		\item The spectrum is mapped onto the Mel scale and energies are aggregated within Mel-spaced frequency bands using overlapping triangular filters.
		\item Mel scale:
		\[
		mel(f) = 1127 \ln \left(1 + \frac{f}{700}\right)
		\]
	\end{itemize}
	\vspace{-6pt}
	\begin{figure}
		\centering
		\hgraphpage[.65\textwidth]{ASR-mel-exp.pdf}
		\caption{Example of a Mel filter bank \cite{2020-jurafsky-martin}}
	\end{figure}
	
\end{frame}

\begin{frame}
	\frametitle{\insertshortsubtitle: \insertsection}
	\framesubtitle{\insertsubsection: Recognition (1)}
	
	\begin{figure}
		\centering
		\hgraphpage[.8\textwidth]{ASR-rec-exp.pdf}
		\caption{Example of speech recognition using an encoder--decoder model \cite{2020-jurafsky-martin}}
	\end{figure}

\end{frame}

\begin{frame}
	\frametitle{\insertshortsubtitle: \insertsection}
	\framesubtitle{\insertsubsection: Recognition (2)}
	
	\begin{itemize}
		\item Sequence probability:
		\[
		P(y_1, \ldots, y_n) = \prod_{i=1}^n P(y_i \mid y_1, \ldots, y_{i-1}, X)
		\]
		
		\item Decoding:
		\[
		\hat{y}_i = \arg\max_{c \in V} P(c \mid y_1, \ldots, y_{i-1}, X)
		\]
		
		\item This behaves as a conditional language model.
		
		\item However, training data may be insufficient to learn a strong language model component.
		
		\item Solution: integrate an external language model.
		
		\[
		score(Y|X) = \frac{1}{|Y|} \log P(Y|X) + \lambda \log P_{LM}(Y)
		\]
	\end{itemize}
	
\end{frame}


\begin{frame}
	\frametitle{\insertshortsubtitle: \insertsection}
	\framesubtitle{\insertsubsection: Some humor}
	
		\begin{center}
			\vgraphpage{humor/humor-ASR.jpg}
		\end{center}
	
\end{frame}

\subsection{Speech synthesis}

\begin{frame}
	\frametitle{\insertshortsubtitle: \insertsection}
	\framesubtitle{\insertsubsection}
	
	\hgraphpage{tts-classif.pdf}
	
\end{frame}

\begin{frame}
	\frametitle{\insertshortsubtitle: \insertsection}
	\framesubtitle{\insertsubsection: Architecture}
	
	\begin{figure}
		\centering
		\hgraphpage[.55\textwidth]{TTS-arch.pdf}
		\caption{Architecture of a voice synthesis system \cite{2017-Hinterleitner}}
	\end{figure}

\end{frame}


\begin{frame}
	\frametitle{\insertshortsubtitle: \insertsection}
	\framesubtitle{\insertsubsection: Some humor}
	
	\begin{center}
		\vgraphpage{humor/humor-TTS.jpg}
	\end{center}

\end{frame}


\insertbibliography{NLP11}{*}

\end{document}

